% v20. General three-high D sign problem: reduction to a nonsingular sector.
% No search bound, parity hypothesis, clutching or node condition is used.
\section{高係数の同時非負性に対する非特異領域への帰着}
\label{hd20:section}

三高分母の完全 matching 定理により，正の weights を持つ整数模型は
\ref{sh14:matching-setup} の canonical 表示で尽くされる．
この節では，三つの高係数が全て非負となる可能性を調べる．
一般の非特異模型についての不存在は主張せず，
その候補から全ての特異 Schur 系と二つの境界を除く．

記号を
\[
 b=-\beta,\quad u=s+t-1,\quad
 A=s\alpha-t,\quad B=-s-tb,\quad C=s+uk,
 \qquad d=U+1-S-T=-D_0
\]
とする．$s,t,S,T,U$ は正整数，$k\ge0$ である．
既知の $\beta=-1$ の境界定理を除いた範囲では $b\ge2$ であり，
$A\ge1$，$C\ge2$ が成り立つ．三つの高係数は
\begin{equation}\label{hd20:D}
 D_a=sd+b-k-1,\qquad
 D_b=td-\alpha+k,\qquad
 D_c=\alpha-b+1-ud
\end{equation}
である．その和は $d$ である．

\begin{theorem}[同時非負性が要求する非特異領域]
\label{hd20:nonsingular}
$q_0=1$，$q_a,q_b,q_c>1$ とし，整数 row sums，pair products，
正の weights を満たす局所模型を考える．
$D_a,D_b,D_c\ge0$ なら，canonical な向き $k\ge0$ において
\[
 b=-\beta\ge2,\qquad d\ge1,\qquad
 \alpha\ge k\ge b+1
\]
であり，Schur 線形系の determinant は
\begin{equation}\label{hd20:positive-det}
 \Delta=C(A-B)+B(A+1)
 \ge(-B)(-B-1)>0
\end{equation}
を満たす．従って，特異な Schur 線形系から三高係数が全て非負である
整数模型が生じることはない．
分母の上限，分母の奇数性，clutching，node 配分，対向 $D$ の相殺は仮定しない．
\end{theorem}

証明を，全係数零の排除と $k=b-1,b$ の境界の排除に分ける．
境界では $d$ の整数性と $S\ge1$ を用いる．

\subsection{全係数零の無上界排除}

\begin{proposition}[三高模型の $D$ ベクトルは零ではない]
\label{hd20:no-zero-D}
上の局所模型において $(D_a,D_b,D_c,D_0)=(0,0,0,0)$ は起こらない．
特に，$D_0=0$ なら三つの高係数には正のものと負のものがともに現れる．
\end{proposition}

\begin{proof}
全係数が零であると仮定する．既知の $\beta=-1$ の境界では $D_0>0$ なので，
$b\ge2$ である．$d=0$ を式~\eqref{hd20:D}へ代入すると
\[
 \alpha=k=b-1
\]
となる．以下
\[
 m=b^2-b+1,\qquad
 K=s^2+st+t^2-s-t,\qquad
 r=(-s,-t,u)^{\mathsf T},\quad x=(-S,-T,U)^{\mathsf T}
\]
と置く．$m$ は三以上の奇整数，$K\ge1$ である．
このとき三次 Schur 補行列は
\[
 B_0=\begin{pmatrix}
 1&b-1&-b\\
 -b&1&b-1\\
 b-1&-b&1
 \end{pmatrix},\qquad
 B_0\mathbf1=B_0^{\mathsf T}\mathbf1=0
\]
である．ここで $B_0$ は三次行列を，$B=-s-tb$ は線形系の scalar を表す．

線形系~\eqref{sh14:matching-linear} の determinant を展開すると
\[
 \Delta=mK>0
\]
となる．$N_S,N_T,N_U$ をそれぞれ $S,T,U$ の Cramer 分子とすると，
直接の多項式恒等式により
\begin{align}
 N_S&=(m-1)\bigl(ms+(B_0r)_a\bigr),\notag\\
 N_T&=(m-1)\bigl(mt+(B_0r)_b\bigr),\notag\\
 N_U&=(m-1)\bigl(mu-(B_0r)_c\bigr),\label{hd20:zero-cramer}\\
 N_U+\Delta-N_S-N_T&=m\bigl(K-(m-1)\bigr).\notag
\end{align}
$U+1-S-T=d=0$ なので最後の式から $K=m-1$ が従う．
残る三式は
\begin{equation}\label{hd20:zero-rotation}
 x=r-\frac1m B_0r
\end{equation}
と書ける．$r,x$ が整数ベクトルであることから，各成分について $m\mid B_0r$ となる．

一方，$J$ を全成分が一の三次行列とすると，直接計算により
\[
 B_0^{\mathsf T}B_0=m(3I-J),\qquad
 \mathbf1^{\mathsf T}r=-1,\qquad
 \lVert r\rVert^2=2K+1=2m-1
\]
である．従って
\[
 \lVert B_0r\rVert^2
 =m\bigl(3\lVert r\rVert^2-1\bigr)=m(6m-4).
\]
左辺は $m^2$ で割り切れるので，$m\mid4$ である．
これは $m\ge3$ が奇整数であることに反する．
最後の主張は $D_a+D_b+D_c=-D_0=0$ と，三つが同時に零とはならないことから従う．
\end{proof}

\subsection{二つの境界の排除}

\begin{lemma}[境界からの同時非負性は生じない]
\label{hd20:k-boundary}
$b\ge2$，$k=b-1$ の canonical 整数模型において
$D_a,D_b,D_c\ge0$ は起こらない．
\end{lemma}

\begin{proof}
非負性と前命題から $d\ge1$ である．$k=b-1$ なら $D_a=sd$ なので，
$D_a\le d$ と $s\ge1$ によって $s=1$ が従う．このとき $u=t$ であり，
Schur 線形系~\eqref{sh14:matching-linear} の第三式は
\[
 t(b-1)T-tbU=b(b-1)>0
\]
となる．ところが $U=S+T+d-1\ge S+T>T$ なので左辺は負となり，矛盾する．
\end{proof}

\begin{lemma}[隣接する境界の排除]
\label{hd20:k-next-boundary}
$b\ge2$，$k=b$ の canonical 整数模型においても
$D_a,D_b,D_c\ge0$ は起こらない．
\end{lemma}

\begin{proof}
前命題により $d\ge1$ である．$D_a=sd-1\le d$ なので
$(s-1)d\le1$ であり，$s=1$ または $(s,d)=(2,1)$ に限られる．
$s=1$ なら線形系の第三式は
\[
 tb(T-U)=b^2-1>0
\]
となるが，やはり $U\ge S+T>T$ なので矛盾する．

$(s,d)=(2,1)$ の場合は $D_a=1,D_b=D_c=0$ であり，
$\alpha=b+t$，$U=S+T$ となる．線形系の第一式と第三式は，それぞれ
\[
 -(2+tb)S+(2b+t)T=b(b+t),
 \qquad bT-(1+tb)S=b^2-1
\]
である．第一式の $b$ 倍から第三式の $(2b+t)$ 倍を引くと
\[
 t(1+b^2+tb)S=t+2b-b^3<t
\]
を得る．$b\ge2,t\ge1$ なので，これは正整数 $S\ge1$ に反する．
両境界の証明に determinant による除算はなく，特異系を落としていない．
\end{proof}

\begin{proof}[定理~\ref{hd20:nonsingular}の証明]
$\beta=-1$ の境界は既知の定理で除かれるので，$b\ge2$ としてよい．
非負性と命題~\ref{hd20:no-zero-D}から $d\ge1$ である．
さらに $0\le D_a\le d$ より
\[
 (s-1)d+b-1\le k\le sd+b-1.
\]
特に $k\ge b-1$ である．補題~\ref{hd20:k-boundary}と
補題~\ref{hd20:k-next-boundary}によって $k=b-1,b$ は起こらず，
整数性から $k\ge b+1$ を得る．また $D_c\ge0$ から
$\alpha\ge ud+b-1$ なので
\[
 \alpha-k\ge(u-s)d=(t-1)d\ge0.
\]
$\bar b=-B=s+tb$ と置くと
\[
 C-\bar b=(s-1)b+u(k-b)\ge0.
\]
従って determinant の恒等式
\[
 \Delta=\bar b(\bar b-1)+(C-\bar b)(A+\bar b)
\]
から式~\eqref{hd20:positive-det}が従う．$\bar b\ge3$ なので右辺は狭義正である．
\end{proof}

\begin{remark}[残る非特異領域の量化]
\label{hd20:remaining}
本定理は全ての特異 Schur 系を，高係数の同時非負性という問題から除くものである．
特異な正整数局所模型自体の不存在を意味しない．
定理~\ref{sh14:singular-family}の無限族はそのまま存在し，
少なくとも一つの高係数が負となる．
また $\Delta>0$ の領域では一般の同時非負性は未解決である．
Cramer 分子の整除性，$S,T,U$ の正整数性，canonical entries の正値性を
引き続き確認する必要がある．
\end{remark}

\begin{example}[残る領域には正整数局所模型がある]
\label{hd20:remaining-example}
$(s,t,\alpha,\beta,k)=(1,2,10,-2,5)$，$(S,T,U)=(4,5,11)$ は
\[
 q=(5,11,23,1),\qquad
 P=\begin{pmatrix}
 0&18&-19&-4\\
 3&0&-9&-5\\
 -6&-28&0&11\\
 -1&-2&2&0
 \end{pmatrix},\qquad
 D=(-1,1,3,-3)
\]
を与える．正の weights は $(5,109/6,45/22,1)$ と取れる．
この例は $k\ge b+1$ や $\Delta>0$ だけを理由に局所模型を除外できないことを示す．
三つの高係数の一つは負であり，同時非負性の反例ではない．
\end{example}

本節の Cramer 恒等式と norm 恒等式は，
\texttt{verify\_high\_d\_boundary\_v20.py} により
多変数整数・有理係数多項式として確認する．
有限範囲の代入から一般定理を推論していない．
